\section{The released \texorpdfstring{$W(E_6)$}{W(E6)} input and action conventions}
\label{sec:released-input}

This section separates the fixed arithmetic premise from the new envelope.
The premise is an exact normal extension $K/\Q$ obtained from the lines of a
smooth cubic surface; the new statements begin only after the exact embedded
subgroups are placed in its labelled action.  This separation matters because
an abstract copy of $\WE$ would not identify any of the fixed fields below.

\subsection{The cubic surface and labelled line action}

Let $Y=V(\Phi)\subset\mathbf P^3_{\Q}$, where
\begin{equation}
\label{eq:surface}
\begin{aligned}
\Phi={}&75081586157u_0^3-28576620789u_0^2u_1-122000922135u_0^2u_2
-5364921951u_0^2u_3\\
&+164150208636u_0u_1^2-415458334296u_0u_1u_2
+151070718312u_0u_1u_3\\
&+1158143874300u_0u_2^2+114691988016u_0u_2u_3
+113572676646u_0u_3^2\\
&+6898957820u_1^3+1132596902196u_1^2u_2-30413540316u_1^2u_3\\
&-2054867641020u_1u_2^2+151980984216u_1u_2u_3
+36794420832u_1u_3^2\\
&+2646295985484u_2^3+560186573940u_2^2u_3
+706181383584u_2u_3^2+1884468968u_3^3.
\end{aligned}
\end{equation}
The released line reconstruction supplies a nonmonic eliminant
\[
 g(d)=L_0d^{27}+b_{26}d^{26}+\cdots+b_0\in\Z[d]
\]
and a fixed ordering $d_1,\ldots,d_{27}$ of its roots corresponding to the
$27$ reconstructed lines.  We reserve $L$ for the degree-$640$ field and use
$L_0>0$ for this leading coefficient.  Set
\begin{equation}
 \alpha_i=L_0d_i\qquad(1\le i\le27).
 \label{eq:scaled-roots}
\end{equation}
The splitting field $K$ of the labelled line carrier is normal and the
faithful line action identifies
\begin{equation}
 G=\Gal(K/\Q)=\WE,\qquad |G|=51840.
 \label{eq:ambient-group}
\end{equation}
The six one-line generators of this action, the line-label bridge, and the
split specialization are recorded in the anonymous exact supplement.  The
surface-to-line realization is consistent with the general cubic-surface
inverse-Galois framework of Elsenhans and Jahnel
\citep{elsenhansjahnel2015moduli}, but the particular coefficients in
\eqref{eq:surface} and their certificate are inputs of this paper, not
attributions to that work.

The normal closure has exact finite ramified support
\begin{equation}
 \begin{gathered}
 \{3,5,181,283,997,1801,2346241,q\},\\
 q=14932047182473291995860108491583652133938007263719.
 \end{gathered}
 \label{eq:released-support}
\end{equation}
At $3$ the released filtration leaves exactly two decomposition-group
possibilities,
\begin{equation}
 D_3=\Tom{140}\quad\text{or}\quad D_3=\Tom{206},
 \qquad I_3=\Tom{140}
 \label{eq:d3-options}
\end{equation}
in the fixed TomLib table.  The locators in \eqref{eq:d3-options} are
versioned checks; the embedded generator arrays and filtration inclusions in
the supplement are the durable data.  We retain both possibilities throughout
\cref{sec:local}.

\subsection{The embedded Gassmann pair}

The durable one-line generators in \cref{app:group-census} define two exact
order-$162$ subgroups $\Hplus,\Hminus\le G$.  Their frozen TomLib locators are
$301$ and $303$, but those integers do not define the subgroups.  Their
indices are $320$, their cores in $G$ are trivial, and their rational
permutation characters are equal.  Thus their fixed fields are arithmetically
equivalent in the sense of the classical Gassmann--Perlis criterion
\citep{perlis1977equation}.  All statements in this paper refer to these
specific embeddings.

For later carrier definitions, let
\begin{equation}
 \cS_+=\Hplus\!\cdot\!\{1,2\}\;\sqcup\;
             \Hplus\!\cdot\!\{1,9\},
 \qquad
 \cS_-=\Hminus\!\cdot\!\{1,2\}.
 \label{eq:basic-supports}
\end{equation}
The two plus orbits have size $27$ each and the minus orbit has size $81$.
The associated integral quadratic sums are
\begin{equation}
 \eta_+=\sum_{\{i,j\}\in\cS_+}\alpha_i\alpha_j,
 \qquad
 \eta_-=\sum_{\{i,j\}\in\cS_-}\alpha_i\alpha_j.
 \label{eq:eta-pm}
\end{equation}
Their formal support stabilizers are $\Hplus$ and $\Hminus$.  The scaling in
\eqref{eq:scaled-roots} is part of the definition; no unscaled sum is used in
an integrality or modular-polynomial claim.

\subsection{Left transport and fixed-field notation}

One-line permutations in the mathematical exposition act on labels on the
left.  The fixed transport is
\begin{equation}
\label{eq:x-array}
x=[1,15,14,13,22,12,27,26,25,7,24,16,17,6,19,18,5,20,4,21,3,2,11,10,9,23,8].
\end{equation}
Thus
\begin{equation}
 \Hthree=x\Hminus x^{-1},\qquad
 \cS_3=x\cS_-,\qquad
 \eta_3=x\eta_-.
 \label{eq:left-transport}
\end{equation}
The same operation is written $\Hminus^x$ in GAP's right-labelled
permutation syntax, because $(i^h)^x=(i^x)^{h^x}$.  This is a conversion of
notation, not an inverse choice.  Exact enumeration gives
\begin{equation}
 \Stab_G(\cS_3)=\Hthree\subset N,
 \label{eq:transport-stabilizer}
\end{equation}
where $N$ is defined intrinsically in \cref{sec:envelope}.  In particular,
$K^{\Hthree}=x(K^{\Hminus})$ is a $\Q$-conjugate of the original minus
field.  We will not identify the two embedded fields.

We use the following standard formulas repeatedly.  If $H\le G$, then the
normal closure of $K^H/\Q$ inside $K$ is fixed by
\[
 \Core_G(H)=\bigcap_{g\in G}gHg^{-1},
\]
and restriction induces
\begin{equation}
 \Aut_{\Q}(K^H)\cong N_G(H)/H.
 \label{eq:normalizer-aut}
\end{equation}
When $J\triangleleft N$, the extension $K^J/K^N$ is Galois with group
$N/J$.  These formulas turn the finite lattice of the next section into the
field-theoretic assertions of \cref{thm:integrated}.
